\section{Complete weak-waveform Fisher operator}
\label{sec:waveform}

The sinusoidal source in Eq.~\eqref{eq:source} is a convenient Fourier probe, but it is not fundamental. The same detector spectrum controls arbitrary finite-dimensional weak temporal waveform estimation within the independent-event model.

Let
\begin{equation}
\Phi_{\boldsymbol\theta}(t)
=\Phi_0\left[1+\sum_{a=1}^{p}\theta_a s_a(t)\right],
\label{eq:waveformSource}
\end{equation}
where the real perturbation modes satisfy $s_a\in L^2(\mathbb R)\cap L^\infty(\mathbb R)$ and $\boldsymbol\theta$ is restricted to a sufficiently small neighborhood of zero so that the intensity remains nonnegative. Define
\begin{equation}
S_a(\omega)=\int_{-\infty}^{\infty}e^{-i\omega t}s_a(t)dt.
\end{equation}
At $\boldsymbol\theta=0$, the incident Poisson Fisher matrix is
\begin{equation}
[F_{\rm in}]_{ab}
=\Phi_0\int s_a(t)s_b(t)dt
=\frac{\Phi_0}{2\pi}
\int S_a^*(\omega)S_b(\omega)d\omega.
\label{eq:waveformInputFI}
\end{equation}

For mark $m$, define the delayed perturbation
\begin{equation}
g_{a,m}(t)
=\int s_a(t-\tau)\mu_m(d\tau).
\end{equation}
The marked output intensity measure is
\begin{equation}
\lambda_{\boldsymbol\theta}(t,dm)
=\Phi_0\left[1+\sum_a\theta_a g_{a,m}(t)\right]\kappa(dm).
\end{equation}

\begin{theorem}[Complete weak-waveform Fisher operator]
For Eq.~\eqref{eq:waveformSource} and the autonomous marked kernel Eq.~\eqref{eq:kernel}, the ideal primary-record Fisher matrix is
\begin{equation}
\boxed{
[F_{\rm out}]_{ab}
=\frac{\Phi_0}{2\pi}
\int_{-\infty}^{\infty}
G(\omega)S_a^*(\omega)S_b(\omega)d\omega,}
\label{eq:waveformFI}
\end{equation}
where $G(\omega)$ is exactly Eq.~\eqref{eq:G}. Parameter-independent downstream processing obeys the matrix data-processing inequality
\begin{equation}
F_{\rm measured}\preceq F_{\rm out}.
\end{equation}
Thus multiplication by $G(\omega)$ is the complete local Fisher-information transfer operator for weak temporal waveform perturbations in this detector class.
\end{theorem}

\begin{proof}
At the baseline,
\begin{equation}
\left.\partial_{\theta_a}\lambda(t,dm)\right|_0
=\Phi_0 g_{a,m}(t)\kappa(dm).
\end{equation}
The Poisson Fisher matrix therefore gives
\begin{equation}
[F_{\rm out}]_{ab}
=\Phi_0\int_{\M}\kappa(dm)
\int g_{a,m}(t)g_{b,m}(t)dt.
\end{equation}
Because
\begin{equation}
\widehat g_{a,m}(\omega)=S_a(\omega)H_m(\omega),
\end{equation}
Parseval's identity yields
\begin{align}
[F_{\rm out}]_{ab}
&=\frac{\Phi_0}{2\pi}
\int S_a^*(\omega)S_b(\omega)
\left[\int_{\M}|H_m(\omega)|^2\kappa(dm)\right]d\omega\\
&=\frac{\Phi_0}{2\pi}
\int G(\omega)S_a^*(\omega)S_b(\omega)d\omega.
\end{align}
The matrix data-processing statement is the usual Fisher-information monotonicity under a parameter-independent stochastic map.
\end{proof}

The long-observation sinusoidal theorem is the Fourier-mode specialization of Eq.~\eqref{eq:waveformFI}: temporal Fourier modes diagonalize the local Fisher operator, with eigenvalue $G(\omega)$. The source choice in Eq.~\eqref{eq:source} is therefore a spectral probe rather than a restriction to a uniquely privileged waveform.

\begin{theorem}[Complete detector ordering]
Consider two detectors $A$ and $B$ in the same theorem class, with transfer spectra $G_A$ and $G_B$. The following are equivalent:
\begin{enumerate}
\item $G_A(\omega)\ge G_B(\omega)$ for almost every $\omega$;
\item for every finite collection of admissible weak temporal perturbations $\{s_a\}$, the corresponding ideal primary-record Fisher matrices obey
\begin{equation}
F_A\succeq F_B.
\end{equation}
\end{enumerate}
\label{thm:ordering}
\end{theorem}

\begin{proof}
If $G_A\ge G_B$ almost everywhere, then for any real coefficient vector $c$, with $S_c=\sum_a c_aS_a$,
\begin{equation}
c^T(F_A-F_B)c
=\frac{\Phi_0}{2\pi}
\int [G_A(\omega)-G_B(\omega)]|S_c(\omega)|^2d\omega
\ge0.
\end{equation}
Hence $F_A-F_B$ is positive semidefinite for every perturbation family.

Conversely, suppose $G_A<G_B$ on a set of nonzero measure. Since $G_A-G_B$ is bounded, measurable, and even, there exists a finite symmetric frequency region containing a positive-measure subset on which $G_A-G_B$ is bounded above by a negative constant. By regularity of Lebesgue measure, choose a real even smooth compactly supported spectrum $S$ concentrated sufficiently strongly on that subset that
\begin{equation}
\int [G_A-G_B]|S|^2d\omega<0.
\end{equation}
Its inverse Fourier transform is a real Schwartz function and hence belongs to $L^2\cap L^\infty$. For this scalar perturbation Eq.~\eqref{eq:waveformFI} gives $F_A<F_B$, contradicting universal Fisher dominance. Thus $G_A\ge G_B$ almost everywhere.
\end{proof}

Theorem~\ref{thm:ordering} gives an operational interpretation of the entire spectrum. If two detectors have the same $G$ almost everywhere, they are Fisher-equivalent for every admissible weak temporal waveform task in this model even if their microscopic mechanisms or latency decompositions differ. If their spectra cross, there is no task-independent ranking: each is superior for some weak temporal perturbations.

\begin{corollary}[Exact band-subspace guarantee]
For a nonzero scalar perturbation $s$ with Fourier transform $S$, define its source-normalized Fisher retention
\begin{equation}
\rho_G[s]
\equiv
\frac{F_{\rm out}[s]}{F_{\rm in}[s]}
=
\frac{\int G(\omega)|S(\omega)|^2d\omega}
{\int |S(\omega)|^2d\omega}.
\label{eq:RayleighG}
\end{equation}
If $S$ is supported in a measurable frequency set $E$, then
\begin{equation}
\operatorname*{ess\,inf}_{\omega\in E}G(\omega)
\le \rho_G[s]\le
\operatorname*{ess\,sup}_{\omega\in E}G(\omega).
\label{eq:bandRayleighBounds}
\end{equation}
Moreover, the infimum and supremum of $\rho_G[s]$ over admissible nonzero perturbations spectrally supported in $E$ equal these two essential bounds.
\end{corollary}

\begin{proof}
Equation~\eqref{eq:RayleighG} is the Rayleigh quotient of the multiplication operator by $G$. The inequalities follow immediately because $|S|^2/\int|S|^2$ is a probability density on frequency. The extremal values are approached by smooth compactly supported spectra concentrated where $G$ approaches its essential infimum or supremum; their inverse Fourier transforms are admissible Schwartz perturbations.
\end{proof}

Therefore preserving at least an absolute Fisher fraction $q$ for \emph{every} admissible weak waveform whose spectrum lies in $|\omega|\le\Omega$ is equivalent to
\begin{equation}
G(\omega)\ge q
\quad\text{for almost every }|\omega|\le\Omega.
\label{eq:universalBandCondition}
\end{equation}
The resource cost of this universal band-subspace guarantee is derived after introducing the integrated timing spectrum below.
